\documentclass[tikz,border=0pt]{standalone}
% Shared TikZ style for figures in
% Layered Network Programming with SNode.C.
%
% Contract for this figure set:
% - every figure PDF has an exact 160 mm canvas width;
% - the TikZ drawing is rendered at native size, not scaled;
% - the drawing is horizontally centered on that canvas;
% - semantic colors, line widths, arrow widths, boundary titles, and spacing
%   are kept consistent unless a local exception is needed to fit the canvas.

\usepackage{xcolor}
\usepackage{tikz}
\usetikzlibrary{
  arrows.meta,
  backgrounds,
  calc,
  fit,
  matrix,
  positioning,
  shapes.geometric,
  shapes.multipart
}

% ----------------------------------------------------------------------
% Palette.
% ----------------------------------------------------------------------
\definecolor{snodecInk}{HTML}{17212B}
\definecolor{snodecMuted}{HTML}{4B5F68}
\definecolor{snodecRule}{HTML}{7A8B93}

\definecolor{snodecBlue}{HTML}{0B4F6C}
\definecolor{snodecBlueSoft}{HTML}{E3F2F8}

\definecolor{snodecGreen}{HTML}{1F7A68}
\definecolor{snodecGreenSoft}{HTML}{E5F4EF}

\definecolor{snodecOrange}{HTML}{C46A1B}
\definecolor{snodecOrangeSoft}{HTML}{FAEFE3}

\definecolor{snodecRed}{HTML}{B24A4A}
\definecolor{snodecRedSoft}{HTML}{F8E8E8}

\definecolor{snodecGraySoft}{HTML}{F1F3F2}

% Title-color palette:
% same hue family as the box, but with stronger contrast.
\colorlet{snodecTitleBlue}{snodecBlue}
\definecolor{snodecTitleGreen}{HTML}{155F51}
\definecolor{snodecTitleOrange}{HTML}{A95516}
\definecolor{snodecTitleGray}{HTML}{4B5F68}
\colorlet{snodecTitleRed}{snodecRed}

% ----------------------------------------------------------------------
% Geometry and line constants.
% ----------------------------------------------------------------------
\newcommand{\snodecCanvasHalfWidth}{80mm}
\newcommand{\snodecCanvasVMargin}{7pt}

\newcommand{\snodecBoxCorner}{2.5pt}
\newcommand{\snodecBoundaryCorner}{3pt}

\newcommand{\snodecBoxLineWidth}{0.70pt}
\newcommand{\snodecInnerLineWidth}{0.58pt}
\newcommand{\snodecBoundaryLineWidth}{0.58pt}
\newcommand{\snodecArrowLineWidth}{0.70pt}
\newcommand{\snodecSecondaryLineWidth}{0.58pt}
\newcommand{\snodecBusLineWidth}{0.58pt}

\newcommand{\snodecBoxXSep}{3mm}
\newcommand{\snodecBoxYSep}{2.5mm}
\newcommand{\snodecSmallBoxXSep}{1.2mm}
\newcommand{\snodecSmallBoxYSep}{1.6mm}
\newcommand{\snodecBoundaryInnerSep}{4.5mm}

\newcommand{\snodecGapSmall}{6mm}
\newcommand{\snodecGapNormal}{8mm}
\newcommand{\snodecGapLarge}{11mm}
\newcommand{\snodecGapMajor}{14mm}
\newcommand{\snodecBendLeg}{6mm}
\newcommand{\snodecOutsideLane}{8mm}

% ----------------------------------------------------------------------
% Shared TikZ styles.
% ----------------------------------------------------------------------
\tikzset{
  snodec diagram/.style={
    font=\sffamily\small,
    text=snodecInk,
    line cap=round,
    line join=round,
    >=Latex
  },
  snodec base box/.style={
    rounded corners=\snodecBoxCorner,
    line width=\snodecBoxLineWidth,
    align=center,
    inner xsep=\snodecBoxXSep,
    inner ysep=\snodecBoxYSep
  },
  snodec runtime box/.style={
    snodec base box,
    draw=snodecBlue!78,
    fill=snodecBlueSoft
  },
  snodec application box/.style={
    snodec base box,
    draw=snodecGreen!78,
    fill=snodecGreenSoft
  },
  snodec protocol box/.style={
    snodec base box,
    draw=snodecOrange!80,
    fill=snodecOrangeSoft
  },
  snodec neutral box/.style={
    snodec base box,
    draw=snodecRule!85,
    fill=snodecGraySoft
  },
  snodec failure box/.style={
    snodec base box,
    draw=snodecRed!82,
    fill=snodecRedSoft
  },
  snodec inner box/.style={
    snodec neutral box,
    line width=\snodecInnerLineWidth,
    inner xsep=\snodecSmallBoxXSep,
    inner ysep=\snodecSmallBoxYSep
  },
  snodec boundary/.style={
    draw=snodecRule!82,
    dashed,
    rounded corners=\snodecBoundaryCorner,
    line width=\snodecBoundaryLineWidth,
    inner sep=\snodecBoundaryInnerSep
  },
  snodec boundary title/.style={
    font=\sffamily\bfseries\small,
    text=snodecTitleBlue,
    fill=white,
    fill opacity=1,
    text opacity=1,
    draw=none,
    anchor=center,
    text height=1.6ex,
    text depth=.25ex,
    inner xsep=1.6pt,
    inner ysep=0.8pt
  },
  snodec arrow label/.style={
    font=\sffamily\bfseries\footnotesize,
    text=snodecInk,
    fill=white,
    fill opacity=1,
    text opacity=1,
    draw=none,
    inner sep=1.0pt
  },
  snodec main arrow/.style={
    -{Latex[length=2.2mm,width=1.6mm]},
    line width=\snodecArrowLineWidth,
    draw=snodecMuted!76
  },
  snodec secondary arrow/.style={
    -{Latex[length=2.0mm,width=1.4mm]},
    line width=\snodecSecondaryLineWidth,
    draw=snodecMuted!68
  },
  snodec bus line/.style={
    draw=snodecMuted!72,
    line width=\snodecBusLineWidth
  },
  snodec return arrow/.style={
    snodec secondary arrow,
    dashed,
    draw=snodecMuted!72
  },
  snodec note/.style={
    font=\sffamily\footnotesize,
    text=snodecMuted,
    align=center
  },
  snodec label/.style={
    font=\sffamily\bfseries\small,
    text=snodecTitleBlue
  },
  % Backward-compatible aliases for old figure sources.
  snodec box/.style={snodec runtime box},
  snodec layer/.style={snodec runtime box,text width=48mm,minimum height=9mm},
  snodec role/.style={snodec application box,text width=38mm},
  snodec protocol/.style={snodec protocol box,text width=42mm},
  snodec neutral/.style={snodec neutral box,text width=42mm},
  snodec failure/.style={snodec failure box,text width=42mm},
  snodec arrow/.style={snodec main arrow,draw=snodecBlue!70},
  snodec muted arrow/.style={snodec secondary arrow}
}

% Fixed-width page/canvas helper. The drawing is not scaled. The natural
% drawing is centered horizontally on a 160 mm wide canvas.
\newcommand{\snodecApplyFixedCanvas}{%
  \coordinate (snodecCanvasCenter) at (current bounding box.center);%
  \coordinate (snodecCanvasSW)%
    at ($(snodecCanvasCenter)+(-\snodecCanvasHalfWidth,0)$ |- current bounding box.south);%
  \coordinate (snodecCanvasNE)%
    at ($(snodecCanvasCenter)+(\snodecCanvasHalfWidth,0)$ |- current bounding box.north);%
  \path[use as bounding box]
    ([yshift=-\snodecCanvasVMargin]snodecCanvasSW)
    rectangle
    ([yshift=\snodecCanvasVMargin]snodecCanvasNE);%
}

\usetikzlibrary{positioning,calc,arrows.meta}

\begin{document}

\def\majorVGap{11mm}
\def\sectionItemGap{3.5mm}
\def\viewGap{5mm}
\def\sectionItemsTopGap{2.6mm}

\begin{tikzpicture}[
  snodec diagram,
  scope block/.style={
    draw=snodecBlue!72,
    fill=snodecBlueSoft,
    rounded corners=\snodecBoxCorner,
    line width=\snodecBoxLineWidth,
    align=center,
    text width=140mm,
    inner xsep=3mm,
    inner ysep=2.4mm,
    font=\sffamily\small
  },
  section block/.style={
    draw=snodecGreen!70,
    fill=snodecGreenSoft,
    rounded corners=\snodecBoxCorner,
    line width=\snodecBoxLineWidth,
    align=center,
    text width=140mm,
    inner xsep=3mm,
    inner ysep=2.4mm,
    font=\sffamily\small
  },
  option block/.style={
    draw=snodecOrange!75,
    fill=snodecOrangeSoft,
    rounded corners=\snodecBoxCorner,
    line width=\snodecBoxLineWidth,
    align=center,
    text width=140mm,
    inner xsep=3mm,
    inner ysep=2.4mm,
    font=\sffamily\small
  },
  section item/.style={
    draw=snodecRule!85,
    fill=snodecGraySoft,
    rounded corners=\snodecBoxCorner,
    line width=\snodecInnerLineWidth,
    align=center,
    text width=21mm,
    inner xsep=1mm,
    inner ysep=1.7mm,
    font=\sffamily\footnotesize
  },
  view block/.style={
    draw=snodecRule!85,
    fill=snodecGraySoft,
    rounded corners=\snodecBoxCorner,
    line width=\snodecInnerLineWidth,
    align=center,
    text width=43mm,
    inner xsep=1mm,
    inner ysep=2.0mm,
    font=\sffamily\footnotesize
  },
  hierarchy flow/.style={
    snodec secondary arrow,
    draw=snodecMuted!75,
    line width=\snodecBoxLineWidth
  }
]

% ----------------------------------------------------------------------
% Main hierarchy.
% ----------------------------------------------------------------------
\node[scope block] (application) {%
  \textbf{\textcolor{snodecTitleBlue}{Application scope}}\\[-0.4mm]
  {\footnotesize executable-level operational shell}\\[-0.2mm]
  {\footnotesize config file, logging, daemonization, help, generated command line}
};

\node[scope block, below=\majorVGap of application] (instance) {%
  \textbf{\textcolor{snodecTitleBlue}{Instance scope}}\\[-0.4mm]
  {\footnotesize configured server/client communication role}\\[-0.2mm]
  {\footnotesize external identity, disabled/enabled state}\\[-0.2mm]
  {\footnotesize e.g. \texttt{echo}, \texttt{mqtt-uplink}, \texttt{measurement-input}}
};

% ----------------------------------------------------------------------
% Section scope.
% ----------------------------------------------------------------------
\node[section block, below=\majorVGap of instance] (section) {%
  \textbf{\textcolor{snodecTitleGreen}{Section scope}}\\[-0.4mm]
  {\footnotesize responsibility boundary inside one instance}\\[\sectionItemsTopGap]
  \begin{tikzpicture}[
    baseline=(connection.base),
    node distance=\sectionItemGap,
    section row item/.style={section item}
  ]
    \node[section row item] (connection) {\texttt{connection}};
    \node[section row item, left=of connection] (remote) {\texttt{remote}};
    \node[section row item, left=of remote] (local) {\texttt{local}};
    \node[section row item, right=of connection] (socket) {\texttt{socket}};
    \node[section row item, right=of socket] (tls) {\texttt{tls}};
  \end{tikzpicture}
};

% ----------------------------------------------------------------------
% Concrete option.
% ----------------------------------------------------------------------
\node[option block, below=\majorVGap of section] (option) {%
  \textbf{\textcolor{snodecTitleOrange}{Option}}\\[-0.4mm]
  {\footnotesize one concrete value inside a section}\\[-0.2mm]
  {\footnotesize host, port, path, timeout, retry, backlog, certificate path}
};

% ----------------------------------------------------------------------
% Hierarchy arrows.
% ----------------------------------------------------------------------
\draw[hierarchy flow] (application.south) -- (instance.north);
\draw[hierarchy flow] (instance.south) -- (section.north);
\draw[hierarchy flow] (section.south) -- (option.north);

% ----------------------------------------------------------------------
% Same hierarchy expressed in three concrete views.
% ----------------------------------------------------------------------
\node[view block, below=\majorVGap of option] (commandLine) {%
  \textbf{\textcolor{snodecTitleGray}{Command line}}\\[-0.4mm]
  \texttt{app instance}\\[-0.2mm]
  \texttt{local --port 8080}
};

\node[view block, left=\viewGap of commandLine] (cppApi) {%
  \textbf{\textcolor{snodecTitleGray}{C++ API}}\\[-0.4mm]
  \texttt{instanceHandle.getConfig()}\\[-0.2mm]
  \texttt{->Local::setPort(8080)}
};

\node[view block, right=\viewGap of commandLine] (configFile) {%
  \textbf{\textcolor{snodecTitleGray}{Config file}}\\[-0.4mm]
  \texttt{instance}\\[-0.2mm]
  \texttt{.local.port = 8080}
};


\snodecApplyFixedCanvas

\end{tikzpicture}
\end{document}
